\section{Data Recipe: los datos no son un inventario, sino un ciclo cerrado}

Solo cuando ya se tienen el robot completo y la cadena de suministro tiene sentido hablar de datos. Este orden no puede invertirse: los datos de robots no son texto descargado de internet, sino algo que se produce en el hardware, los escenarios, las tareas y la retroalimentación de los clientes.

El «panorama de datos» de EP132 y el de EP134 se aclaran mutuamente. Lo central del Data Recipe del que habla aquí Gao Jiyang (高继扬) no es vender los datos como producto, sino conectar el robot completo, los escenarios, las tareas, la recolección, el entrenamiento, la evaluación y el ecosistema de código abierto en un volante iterable.

\begin{figure}[H]
\centering
\includegraphics[width=0.96\textwidth]{robot-data-recipe-flywheel.png}
\caption{Volante del Data Recipe para robots: despliegue del robot completo, recolección de datos, entrenamiento del modelo, valor en campo, retorno de errores y apertura compartida.}
\end{figure}

\begin{knowledgebox}{Lectura de la figura: por qué el Data Recipe no consiste en vender datos}
Los datos de la figura provienen de hardware y escenarios reales. Cuando el robot completo llega a escenarios de desarrolladores o de productividad, genera trayectorias, fallas y retroalimentación; estos datos entrenan VLA, modelos base y herramientas de posentrenamiento; cuando el modelo crea valor para el cliente, trae consigo más escenarios y más retroalimentación. Los datos no son un inventario estático, sino el combustible de aprendizaje que generan de manera continua el producto y los clientes.
\end{knowledgebox}

\subsection{Por qué importa llegar a diez mil unidades entregadas}

Gao Jiyang considera que «lograr entregas de diez mil unidades en escenarios de productividad» es una bet importante en este momento. Esta meta no es una simple cifra de ventas, sino una validación integral de los datos, la cadena de suministro, la calidad y el valor para el cliente. Diez mil unidades significan que el hardware puede entregarse, que los clientes están dispuestos a usarlo, que el sistema de mantenimiento puede seguirle el paso y que el retorno de datos puede escalarse.

\begin{knowledgebox}{Términos clave: el ciclo cerrado de datos de robots}
\begin{tabularx}{\textwidth}{p{0.23\textwidth}p{0.32\textwidth}X}
\toprule
Término & Problema que resuelve & Significado en este episodio \\
\midrule
Data Recipe & Receta de origen, selección, entrenamiento y retroalimentación de los datos & Determina cómo los datos reales se convierten en capacidades del modelo. \\
Demo in Video & Demostración que parece funcionar en un video & Tiende a generar expectativas demasiado altas. \\
Demo in Office & Demostración en vivo en la oficina de la empresa & Más sólida que el video, pero sigue siendo un entorno controlado. \\
Demo in the Wild & Demostración desplegada con clientes, en ferias y en escenarios de distintos países & Más cercana a la capacidad real de generalización. \\
Escenario de productividad & Los clientes usan el robot para realizar trabajo real & Fuente del ciclo cerrado entre datos y valor comercial. \\
\bottomrule
\end{tabularx}
\end{knowledgebox}

\subsection{Código abierto y datos compartidos}

Xinghaitu (星海图) publica en código abierto conjuntos de datos y modelos, y comparte datos con sus clientes. Gao Jiyang subraya que los datos no son su negocio principal, pero que con gusto los comparte. Esto puede entenderse como una estrategia de ecosistema: permitir que más desarrolladores y clientes validen tareas, hagan visibles sus necesidades y generen retroalimentación, para así ampliar los escenarios de uso de los modelos base y de la cadena de herramientas.

\subsection{Resumen de la sección}

La esencia del Data Recipe es conectar el robot completo, los clientes y el entrenamiento de modelos. Lo que de verdad escasea en una empresa de robótica no es «tener un lote de datos», sino un mecanismo que produzca de forma continua datos, evaluaciones y retroalimentación de alto valor.
